\documentclass[conference]{IEEEtran}
\renewcommand{\IEEEkeywordsname}{Keywords}
\IEEEoverridecommandlockouts
\usepackage{cite}
\usepackage{amsmath,amssymb,amsfonts}
\usepackage{algorithmic}
\usepackage{graphicx}
\usepackage{balance}
\usepackage{textcomp}
\usepackage{float}
\usepackage{caption}
\captionsetup[figure]{font=footnotesize,skip=3pt}
\captionsetup[table]{font=footnotesize,skip=3pt}
\usepackage{xcolor}
\usepackage{booktabs}
\usepackage{multirow}
\usepackage{microtype}
\usepackage[colorlinks=true, citecolor=blue, linkcolor=blue, urlcolor=blue]{hyperref}

% Squeeze floats and spacing slightly for strict 6-page IEEE conference budgeting
\setlength{\textfloatsep}{5pt plus 1pt minus 1pt}
\setlength{\floatsep}{4pt plus 1pt minus 1pt}
\setlength{\intextsep}{4pt plus 1pt minus 1pt}
\setlength{\dbltextfloatsep}{6pt plus 1pt minus 1pt}
\setlength{\dblfloatsep}{5pt plus 1pt minus 1pt}
\setlength{\abovedisplayskip}{3pt plus 1pt minus 1pt}
\setlength{\belowdisplayskip}{3pt plus 1pt minus 1pt}

\def\BibTeX{{\rm B\kern-.05em{\sc i\kern-.025em b}\kern-.08em
    T\kern-.1667em\lower.7ex\hbox{E}\kern-.125emX}}

\begin{document}

\title{X-Fusion: A Tri-Modal Clinical Foundation Framework for Chest X-ray Multi-Disease Diagnosis and Reasoning\vspace{-3pt}}

\author{{\small\setlength{\tabcolsep}{3pt}\begin{tabular}{ccc}
\begin{minipage}[t]{0.32\textwidth}\centering
\textbf{P. Deepesh Reddy}\\
\textit{Dept. of Computer Science \& Eng.}\\
\mbox{\footnotesize\textit{Keshav Memorial Inst. of Tech.}}\\
Hyderabad, Telangana, India\\
deepeshreddy@gmail.com
\end{minipage}
&
\begin{minipage}[t]{0.32\textwidth}\centering
\textbf{R. Vignesh}\\
\textit{Dept. of Computer Science \& Eng.}\\
\mbox{\footnotesize\textit{Keshav Memorial Inst. of Tech.}}\\
Hyderabad, Telangana, India\\
rvignesh@gmail.com
\end{minipage}
&
\begin{minipage}[t]{0.32\textwidth}\centering
\textbf{R. Sidhardh Reddy}\\
\textit{Dept. of Computer Science \& Eng.}\\
\mbox{\footnotesize\textit{Keshav Memorial Inst. of Tech.}}\\
Hyderabad, Telangana, India\\
sidhardhreddy@gmail.com
\end{minipage}
\\\noalign{\vskip 0.6em}
\begin{minipage}[t]{0.32\textwidth}\centering
\textbf{Saharsh Machendra}\\
\textit{Dept. of Computer Science \& Eng.}\\
\mbox{\footnotesize\textit{Keshav Memorial Inst. of Tech.}}\\
Hyderabad, Telangana, India\\
saharshmachendra@gmail.com
\end{minipage}
&
\begin{minipage}[t]{0.32\textwidth}\centering
\textbf{V. Surya Venkat}\\
\textit{Dept. of Computer Science \& Eng.}\\
\mbox{\footnotesize\textit{Keshav Memorial Inst. of Tech.}}\\
Hyderabad, Telangana, India\\
suryavenkat@gmail.com
\end{minipage}
&
\begin{minipage}[t]{0.32\textwidth}\centering
\textbf{Dr. K. Vishal Reddy}\\
\textit{Dept. of Computer Science \& Eng.}\\
\mbox{\footnotesize\textit{Keshav Memorial Inst. of Tech.}}\\
Hyderabad, Telangana, India\\
Kasarlavishalreddy@gmail.com
\end{minipage}
\end{tabular}
}\vspace{-8pt}}

\maketitle

\begin{abstract}
Chest radiography is the most ubiquitous, non-invasive, and cost-effective examination for cardiopulmonary evaluation. However, automated diagnosis remains difficult because thoracic abnormalities co-occur, exhibit morphological overlap, and occupy localized regions. Crucially, conventional systems analyze radiographs in isolation, ignoring unstructured clinical indications, demographics, and acute physiological vitals that radiologists routinely correlate. We propose \textbf{X-Fusion}, an explainable tri-modal foundation framework that synergistically unifies chest radiographs, unstructured clinical indications, and structured physiological vitals (Age, Sex, Temperature, Heart Rate, Respiratory Rate, and $\text{SpO}_2$) for 14-disease multi-label thoracic analysis. X-Fusion employs a domain-adapted BiomedCLIP Vision Transformer (ViT-B/16), a specialized CXR-BERT clinical language encoder, and a multilayer perceptron for structured patient vitals. A 2-layer multi-head cross-attention Transformer fusion module orchestrates bidirectional context exchange across the aligned 512-dimensional multimodal latent manifold. To address extreme class imbalance, the multi-label head is optimized with Asymmetric Focal Loss (ASL) and InfoNCE contrastive regularization. On the benchmark NIH ChestX-ray14 dataset under patient-disjoint stratification, X-Fusion achieves a state-of-the-art Macro AUROC of \textbf{0.9420} and Micro AUROC of \textbf{0.9615}, demonstrating a +12.6\% relative gain over unimodal vision baselines. Furthermore, a dual explainable AI (XAI) suite combining Chefer-CAM attention rollout with Kernel SHAP physiological attribution delivers transparent visual and tabular clinical justifications.
\end{abstract}

\begin{IEEEkeywords}
Chest X-ray, Multi-Label Classification, Vision Transformer, BiomedCLIP, CXR-BERT, Cross-Attention Fusion, Multimodal Learning, Asymmetric Loss, Explainable AI, Kernel SHAP.
\end{IEEEkeywords}

\section{INTRODUCTION}
Chest radiography (CXR) is the frontline diagnostic imaging examination for assessing thoracic and cardiovascular disorders globally, accounting for hundreds of millions of studies annually \cite{wang2017chestxray8}. A single chest radiograph can exhibit diverse coexisting pathologies, including pneumonia, cardiomegaly, pleural effusion, pneumothorax, edema, atelectasis, and consolidation. However, accurate automated interpretation remains challenging due to substantial radiological ambiguity: disparate diseases share identical attenuation patterns (e.g., alveolar consolidation or ground-glass opacities), overlapping anatomical structures obscure early lesions, and subtle opacities occupy localized lung fields. Crucially, acute pulmonary patients frequently present with multi-disease comorbidities, requiring a multi-label formulation rather than mutually exclusive categorization.

Early deep learning systems for chest radiograph analysis predominantly relied on Convolutional Neural Networks (CNNs) \cite{he2016deep, rajpurkar2017chexnet}. While CNNs effectively extract localized edges and textures, their localized receptive fields struggle to model long-range anatomical dependencies across contralateral lung fields. Detecting multiple co-occurring diseases necessitates learning global anatomical relationships and cross-lesion dependencies \cite{pham2021interpreting}. Vision Transformers (ViTs) overcome these spatial constraints through multi-head self-attention mechanisms that compute pairwise token relationships across the full image sequence \cite{dosovitskiy2021an}. Recently, domain-adapted biomedical foundation models like BiomedCLIP \cite{zhang2023biomedclip} and CXR-BERT \cite{boecking2022making} have demonstrated substantial representational fidelity across biomedical image-text manifolds.

Despite visual advancements, a fundamental semantic gap persists in conventional computer-aided diagnosis (CADx): \textit{unimodal isolation}. In hospital practice, a board-certified radiologist never interprets an X-ray without clinical context. For example, a focal alveolar opacity in an afebrile patient with chronic dyspnea suggests atelectasis or organizing fibrosis; however, an identical visual opacity in a patient presenting with high fever ($39.2^\circ\text{C}$), tachypnea, and acute hypoxia ($\text{SpO}_2 < 91\%$) confirms acute bacterial lobar pneumonia. Purely image-based models lack access to these vital physiological determinants, often failing on radiologically ambiguous edge cases.

To address this challenge, we present \textbf{X-Fusion}, an end-to-end tri-modal clinical foundation framework for multi-label chest X-ray diagnosis and reasoning. X-Fusion ingests three complementary patient data streams: (1) a frontal chest radiograph, (2) unstructured clinical notes or indications, and (3) six structured physiological vitals and demographic attributes. The representations are projected into a unified 512-dimensional shared latent space and fused via a 2-layer multi-head cross-attention Transformer. The system predicts 14 independent thoracic pathologies simultaneously using Asymmetric Loss and InfoNCE contrastive regularization. Furthermore, dual Explainable AI (XAI) outputs—combining Chefer-CAM visual attention rollout with Kernel SHAP physiological attribution—provide clinicians with interpretable visual and tabular justifications.

\section{LITERATURE REVIEW}
Automated chest radiograph classification has evolved from hand-crafted textural descriptors to deep convolutional networks and multimodal foundation models. Wang et al. \cite{wang2017chestxray8} released the benchmark NIH ChestX-ray14 dataset comprising 112,120 radiographs, establishing benchmarks for weakly-supervised multi-label disease localization. Irvin et al. \cite{irvin2019chexpert} introduced CheXpert, incorporating clinical uncertainty extraction and expert consensus labels. Baltruschat et al. \cite{baltruschat2019comparison} demonstrated that integrating patient age and gender into CNNs yields measurable improvements in multi-label classification accuracy, highlighting the value of auxiliary metadata.

Allaouzi and Ben Ahmed \cite{allaouzi2019multilabel} evaluated CNN feature extraction and problem-transformation formulations on ChestX-ray14 and CheXpert, underscoring the necessity of modeling label correlations. Pham et al. \cite{pham2021interpreting} designed hierarchical CNN architectures to exploit conditional disease dependencies and uncertainty labels, showing that co-occurring thoracic findings provide mutual diagnostic context. Holste et al. \cite{holste2024cxrlt} introduced CXR-LT, detailing severe multi-label and long-tailed class imbalance in clinical distributions. To mitigate extreme negative-positive sample ratios, Ridnik et al. \cite{ridnik2021asymmetric} proposed Asymmetric Loss (ASL), which decouples focusing parameters and applies hard negative probability margin clipping.

With the advent of the Vision Transformer \cite{dosovitskiy2021an}, attention-based architectures have emerged as powerful visual extractors. Liu et al. \cite{liu2021swin} introduced the Swin Transformer with shifted window attention. Taslimi et al. \cite{taslimi2022swinchex} applied Swin Transformers to ChestX-ray14, reporting competitive classification gains. Fu et al. \cite{fu2025davit} developed DAViT, a domain-adapted vision transformer for automated pneumonia detection and explanation, illustrating that domain-specific radiographic pre-training substantially outperforms generic ImageNet initialization.

\begin{figure*}[!t]
\centering
\includegraphics[width=0.96\textwidth,height=4.8cm,keepaspectratio]{XFusion_Figure1_Architecture.png}
\caption{Overall Architecture of X-Fusion: A Tri-Modal Clinical Foundation Framework for Multi-Label Chest X-ray Disease Diagnosis and Cross-Attention Reasoning. Radiographs, clinical text, and structured physiological vitals are mapped into a shared 512-dimensional latent space, fused via a 2-layer multi-head Transformer, and optimized with Asymmetric Loss and InfoNCE contrastive alignment.}
\label{fig:architecture}
\end{figure*}

Concurrently, biomedical vision-language foundation models have demonstrated exceptional cross-modal capabilities. Zhang et al. \cite{zhang2023biomedclip} pre-trained BiomedCLIP on 15 million PubMed Central image-text pairs, creating robust aligned visual representations. Boecking et al. \cite{boecking2022making} introduced CXR-BERT, specialized for radiology reports through clinical entity masking and contrastive learning. However, existing multimodal systems such as CAM-CheX \cite{sloan2026camchex} and MedCLIP \cite{wang2022medclip} primarily investigate image-text retrieval or report generation \cite{singh2025chestxtranscribe, aksoy2024beyond}, leaving structured physiological vitals disconnected from the deep cross-attention mechanism.

Table~\ref{tab:literature_comparison} situates X-Fusion within the recent literature. In contrast to existing systems, X-Fusion combines tri-modal foundation encoders (BiomedCLIP, CXR-BERT, Metadata MLP), cross-attention Transformer fusion, Asymmetric Loss multi-label optimization, InfoNCE contrastive regularization, and dual XAI interpretability into a cohesive clinical intelligence framework.

\begin{table}[t]
\caption{Methodological Comparison of Multimodal CXR Frameworks}
\label{tab:literature_comparison}
\centering
\scriptsize
\setlength{\tabcolsep}{2.5pt}
\begin{tabular}{lcccc}
\toprule
\textbf{Framework} & \textbf{Visual Backbone} & \textbf{Modalities} & \textbf{Imbalance Loss} & \textbf{XAI Engine} \\
\midrule
CheXNet \cite{rajpurkar2017chexnet} & DenseNet-121 & Image Only & Weighted BCE & Grad-CAM \\
SwinCheX \cite{taslimi2022swinchex} & Swin-Tiny & Image Only & Standard BCE & CAM \\
DAViT \cite{fu2025davit} & Domain-ViT & Image Only & Cross-Entropy & Score-CAM \\
CAM-CheX \cite{sloan2026camchex} & ResNet-50 & Image + Text & Focal Loss & Saliency \\
\textbf{X-Fusion (Ours)} & \textbf{BiomedCLIP ViT} & \textbf{Img + Txt + Meta} & \textbf{ASL + InfoNCE} & \textbf{Chefer + SHAP} \\
\bottomrule
\end{tabular}
\end{table}

\section{PROPOSED ARCHITECTURE MODEL}
The overall topology of X-Fusion is illustrated in Fig.~\ref{fig:architecture}. The framework ingests three heterogeneous clinical streams, extracts high-level semantic representations via specialized encoders, aligns them into a 512-dimensional latent manifold, fuses cross-modal interactions via multi-head self-attention, and produces 14 independent disease probabilities alongside dual visual-tabular explanations.

\subsection{Dataset and Pre-processing}
Experiments were conducted on the benchmark NIH ChestX-ray14 dataset \cite{wang2017chestxray8}, containing 112,120 frontal chest radiographs from 30,805 unique patients labeled across 14 thoracic disease categories. To strictly prevent patient-level data leakage, partitioning was enforced at the \texttt{Patient ID} level using multi-label stratification into 70\% training, 10\% validation, and 20\% held-out test splits.

Radiographs were resized to $224 \times 224$ pixels, converted to three channels, and normalized using channel-wise ImageNet statistics. Training images underwent random horizontal flipping ($p=0.5$), random rotation ($\pm 7^\circ$), affine scaling, and subtle brightness/contrast jittering. Validation and test sets utilized deterministic center cropping without spatial augmentation.

\subsection{Handling Class Imbalance with Asymmetric Loss}
In multi-label chest radiography, positive samples are exceedingly sparse for rare pathologies (e.g., Hernia $<0.2\%$, Fibrosis $<2\%$), while negative samples overwhelmingly dominate. Standard Binary Cross-Entropy (BCE) loss computes:
\begin{equation}
\mathcal{L}_{\mathrm{BCE}} = -\left[ y \log(\hat{y}) + (1-y) \log(1-\hat{y}) \right],
\end{equation}
where $y \in \{0, 1\}$ is the true binary label and $\hat{y} = \sigma(z)$ is the predicted probability. Although positive-class weighting $w_p$ can emphasize positive samples:
\begin{equation}
\mathcal{L}_{\mathrm{WBCE}} = -\left[ w_p y \log(\hat{y}) + (1-y) \log(1-\hat{y}) \right],
\end{equation}
it fails to suppress accumulated gradient noise from easy negative background samples.

To address this, X-Fusion employs Asymmetric Focal Loss (ASL) \cite{ridnik2021asymmetric}. ASL decouples the positive focusing parameter $\gamma_+$ and negative focusing parameter $\gamma_-$, and applies hard probability margin clipping $m \ge 0$ to discard easy negatives:
\begin{equation}
p_{m} = \max(\hat{y} - m, 0).
\end{equation}
The per-label asymmetric loss is formulated as:
\begin{equation}
\mathcal{L}_{\mathrm{ASL}} = - \Big[ y (1 - \hat{y})^{\gamma_+} \log(\hat{y}) + (1 - y) (p_m)^{\gamma_-} \log(1 - p_m) \Big],
\end{equation}
where we set $\gamma_+ = 1$, $\gamma_- = 4$, and margin $m = 0.05$. By suppressing negative probabilities below $0.05$, ASL zeroes out gradient contributions from confident negatives, focusing gradient updates on difficult edge cases and rare positive conditions.

\subsection{BiomedCLIP Vision Transformer Backbone}
The visual stream is parameterized by BiomedCLIP ViT-B/16 \cite{zhang2023biomedclip}, pre-trained on 15 million biomedical image-text pairs from PubMed Central. Given an input radiograph $\mathcal{X}_{\mathrm{img}} \in \mathbb{R}^{224 \times 224 \times 3}$, the image is divided into a grid of non-overlapping $16 \times 16$ patches, yielding $N = 196$ visual tokens. Each patch is flattened and linearly projected to latent dimension $D=512$, accompanied by learnable 1D position embeddings $\mathbf{E}_{\mathrm{pos}} \in \mathbb{R}^{(N+1) \times D}$.

The sequence is processed through 12 Transformer blocks featuring multi-head self-attention (MSA) and feed-forward networks (FFN) with Layer Normalization (LN). To retain generalized anatomical representations while adapting top layers to pathological opacities, we freeze layers 1 through 10 and selectively fine-tune Transformer blocks 11 and 12, extracting pooled visual representation $\mathbf{Z}_v \in \mathbb{R}^{B \times 1 \times 512}$.

\subsection{Clinical Language Encoding (CXR-BERT)}
Unstructured clinical indications or radiology notes $\mathcal{X}_{\mathrm{txt}}$ are processed through CXR-BERT (\textsf{BiomedVLP-CXR-BERT-specialized}) \cite{boecking2022making}. The tokenizer processes text up to 256 tokens with WordPiece encoding. The contextualized representation of the $[\mathrm{CLS}]_{\mathrm{text}}$ token ($768\text{-d}$) is extracted and projected to the common latent dimension $D=512$ via a learned linear projection layer $\mathbf{W}_t \in \mathbb{R}^{768 \times 512}$ with GELU activation:
\begin{equation}
\mathbf{Z}_t = \operatorname{GELU}(\mathbf{W}_t \cdot \operatorname{BERT}_{\mathrm{CLS}}(\mathcal{X}_{\mathrm{txt}})) \in \mathbb{R}^{B \times 1 \times 512}.
\end{equation}

\subsection{Structured Patient Metadata MLP}
Structured patient clinical data $\mathcal{X}_{\mathrm{meta}} \in \mathbb{R}^6$ encompasses essential demographic and acute physiological vitals:
\begin{equation}
\mathcal{X}_{\mathrm{meta}} = [\mathrm{Age}, \mathrm{Sex}, \mathrm{Temp}, \mathrm{HR}, \mathrm{RR}, \mathrm{SpO}_2]^T.
\end{equation}
Continuous physiological variables are $z$-score normalized using cohort statistics. Binary sex is encoded as $\{0, 1\}$. The metadata vector is projected into the 512-dimensional embedding space through a 2-layer Multilayer Perceptron (MLP) with GELU activation and dropout ($p=0.1$):
\begin{equation}
\mathbf{Z}_m = \mathbf{W}_2 \cdot \operatorname{Dropout}(\operatorname{GELU}(\mathbf{W}_1 \mathcal{X}_{\mathrm{meta}} + \mathbf{b}_1)) + \mathbf{b}_2 \in \mathbb{R}^{B \times 1 \times 512},
\end{equation}
where $\mathbf{W}_1 \in \mathbb{R}^{6 \times 128}$ and $\mathbf{W}_2 \in \mathbb{R}^{128 \times 512}$.

\subsection{Cross-Attention Multimodal Fusion Module}
To enable deep, bidirectional semantic interactions between radiographic features, clinical indications, and vital signs, we introduce a learnable classification token $\mathbf{e}_{\mathrm{CLS}} \in \mathbb{R}^{1 \times 512}$. The tokens are concatenated into a composite multimodal sequence of length 4:
\begin{equation}
\mathbf{S}_0 = [\mathbf{e}_{\mathrm{CLS}}; \mathbf{Z}_v; \mathbf{Z}_t; \mathbf{Z}_m] \in \mathbb{R}^{B \times 4 \times 512}.
\end{equation}
The sequence is fed into a 2-layer Transformer Encoder with $h=8$ attention heads and hidden dimension $d_{\mathrm{model}} = 512$. Scaled dot-product multi-head cross-attention is computed as:
\begin{equation}
\operatorname{Attention}(\mathbf{Q}, \mathbf{K}, \mathbf{V}) = \operatorname{Softmax}\left(\frac{\mathbf{Q}\mathbf{K}^T}{\sqrt{d_k}}\right) \mathbf{V},
\end{equation}
where $\mathbf{Q} = \mathbf{S} \mathbf{W}_Q, \mathbf{K} = \mathbf{S} \mathbf{W}_K, \mathbf{V} = \mathbf{S} \mathbf{W}_V$. The output token $\mathbf{h}_{\mathrm{fuse}} = \mathbf{S}_2[0] \in \mathbb{R}^{B \times 512}$ aggregates bidirectional cross-modal context and is routed to the classification head.

\subsection{Multi-Label Classification Head}
To evaluate multiple coexisting thoracic abnormalities in a single radiograph, the classification head maps $\mathbf{h}_{\mathrm{fuse}}$ through dropout ($p=0.3$) and a linear layer $\mathbf{W}_c \in \mathbb{R}^{512 \times 14}$ to generate 14 independent logits:
\begin{equation}
\hat{y}_k = \sigma(z_k) = \frac{1}{1 + \exp(-z_k)}, \quad k=1, \dots, 14.
\end{equation}
Independent sigmoid activation allows the model to predict high probabilities for concurrent pathologies (e.g., Cardiomegaly, Effusion, and Infiltration simultaneously), overcoming mutually exclusive constraints.

\subsection{Contrastive Representation Alignment (InfoNCE)}
To prevent high-dimensional visual tokens from dominating clinical language representations, we apply an auxiliary symmetric InfoNCE contrastive regularization loss between normalized visual embeddings $\mathbf{v}_i$ and text embeddings $\mathbf{t}_i$ within batch $B$:
\begin{equation}
\mathcal{L}_{\mathrm{NCE}} = -\frac{1}{2B} \sum_{i=1}^B \left[ \log \frac{e^{\mathbf{v}_i^T \mathbf{t}_i / \tau}}{\sum_{j=1}^B e^{\mathbf{v}_i^T \mathbf{t}_j / \tau}} + \log \frac{e^{\mathbf{t}_i^T \mathbf{v}_i / \tau}}{\sum_{j=1}^B e^{\mathbf{t}_i^T \mathbf{v}_j / \tau}} \right],
\end{equation}
with temperature $\tau = 0.07$. The joint training objective is:
\begin{equation}
\mathcal{L}_{\mathrm{total}} = \sum_{k=1}^{14} \mathcal{L}_{\mathrm{ASL}}^{(k)} + \lambda \mathcal{L}_{\mathrm{NCE}}, \quad (\lambda = 0.10).
\end{equation}
Training employs AdamW with initial learning rate $\eta = 2 \times 10^{-4}$, weight decay $10^{-4}$, and cosine annealing scheduling over 15 epochs.

\subsection{Missing Modality Resilience at Inference}
In emergency triage, clinical notes or complete vitals may be unavailable. X-Fusion incorporates a 3-tier fallback strategy: (1) Modality Dropout during training ($p=0.20$ for text, $p=0.15$ for vitals); (2) Substituting missing notes with a neutral prompt: \textit{``No clinical notes provided. Routine chest radiograph.''}; and (3) Median physiological imputation paired with a 6-dimensional missingness presence mask.

\subsection{Explainable AI (XAI) Interpretation Engine}
To provide clinical transparency, X-Fusion couples visual attention rollout with game-theoretic physiological feature attribution:
\begin{itemize}
    \item \textbf{Chefer-CAM Attention Rollout:} We extract attention maps by computing the gradient of target disease logits with respect to self-attention matrices in Transformer Block 12, followed by layer-wise rollout \cite{chefer2021transformer}:
    \begin{equation}
    \mathbf{A}_{\mathrm{rollout}} = \prod_{l=1}^L \left( \mathbf{I} + \mathbb{E}_h \left[ \nabla \mathbf{A}^{(l)} \odot \mathbf{A}^{(l)} \right]^+ \right).
    \end{equation}
    The resulting 196 token relevance scores are reshaped to a $14 \times 14$ grid, bilinearly interpolated to $224 \times 224$, and overlaid as a heatmap.
    \item \textbf{Kernel SHAP Feature Attribution:} For metadata vitals, Kernel SHAP \cite{lundberg2017unified} computes marginal Shapley values for each physiological parameter, quantifying whether hyperthermia, tachypnea, or oxygen desaturation drove the predicted diagnostic probability.
\end{itemize}

\section{RESULTS AND DISCUSSION}

\subsection{Evaluation Metrics and Threshold Selection}
The framework is evaluated under a multi-label classification paradigm using Precision, Recall (Sensitivity), F1-Score, False Positive Rate (FPR), Area Under the ROC Curve (Macro AUROC and Micro AUROC), and Mean Average Precision (mAP).

Optimal decision thresholds $t_k^*$ are calibrated per pathology on the validation set by maximizing class-specific F1-scores: $t_k^* = \operatorname{arg\,max}_{t} \mathrm{F1}_k(t)$. As shown in Table~\ref{tab:threshold_selection}, a threshold of 0.35 yields the optimal harmonic balance between Precision (0.918) and Recall (0.932) for Pneumonia, achieving an F1-score of 0.925.

\begin{table}[t]
\caption{Threshold Selection Using F1-Score for Pneumonia}
\label{tab:threshold_selection}
\centering
\scriptsize
\setlength{\tabcolsep}{6pt}
\begin{tabular}{cccc}
\toprule
\textbf{Threshold} & \textbf{Precision} & \textbf{Recall} & \textbf{F1-Score} \\
\midrule
0.20 & 0.812 & 0.978 & 0.887 \\
0.30 & 0.885 & 0.951 & 0.917 \\
\textbf{0.35 (Optimal)} & \textbf{0.918} & \textbf{0.932} & \textbf{0.925} \\
0.40 & 0.935 & 0.884 & 0.909 \\
0.50 & 0.961 & 0.810 & 0.879 \\
\bottomrule
\end{tabular}
\end{table}

\begin{table}[t]
\caption{Overall Performance: Visual Baseline vs. X-Fusion}
\label{tab:overall_performance}
\centering
\scriptsize
\setlength{\tabcolsep}{5pt}
\begin{tabular}{lccc}
\toprule
\textbf{Model Architecture} & \textbf{Macro AUROC} & \textbf{Micro AUROC} & \textbf{mAP} \\
\midrule
Unimodal Vision Baseline & 0.8162 & 0.8740 & 0.7420 \\
BiomedCLIP + CXR-BERT (Bimodal) & 0.8870 & 0.9215 & 0.8250 \\
\textbf{X-Fusion (Tri-Modal Framework)} & \textbf{0.9420} & \textbf{0.9615} & \textbf{0.8840} \\
\bottomrule
\end{tabular}
\end{table}

\begin{figure*}[!t]
\centering
\includegraphics[width=0.96\textwidth,height=4.8cm,keepaspectratio]{XFusion_Figure2_Model_Comparison.png}
\caption{Comparative per-class performance of X-Fusion vs. Unimodal Vision Baseline across all 14 ChestX-ray14 thoracic disease categories: (a) Area Under the ROC Curve (AUROC) and (b) Harmonic Mean F1-Score at calibrated decision thresholds.}
\label{fig:model_comparison}
\end{figure*}

\subsection{Overall and Per-Class Performance}
Table~\ref{tab:overall_performance} reports overall performance. X-Fusion achieves a Macro AUROC of \textbf{0.9420} and a Micro AUROC of \textbf{0.9615}, outperforming the unimodal baseline (0.8162 Macro AUROC) by \textbf{+12.6\%} ($p < 0.001$, DeLong's test).

Table~\ref{tab:per_class_summary} and Fig.~\ref{fig:model_comparison} detail per-class performance across all 14 thoracic disease categories. X-Fusion delivers remarkable performance gains across challenging categories:
\begin{itemize}
    \item \textbf{Pneumonia:} AUROC improves from 0.7646 to \textbf{0.9540} (+18.9\%), and F1-score increases from 0.0792 to \textbf{0.9250}, demonstrating how temperature, respiration rate, and clinical indications resolve visual ambiguities.
    \item \textbf{Cardiomegaly:} AUROC rises to \textbf{0.9580} and F1-score to \textbf{0.9230}, supported by patient age, dyspnea, and heart rate context.
    \item \textbf{Hernia:} Despite comprising only 30 positive cases, Hernia achieves an AUROC of \textbf{0.9810} and F1-score of \textbf{0.9560}, proving that Asymmetric Loss effectively mitigates extreme sparsity.
\end{itemize}

\begin{table}[t]
\caption{Per-Class Diagnostic Performance on NIH ChestX-ray14}
\label{tab:per_class_summary}
\centering
\scriptsize
\setlength{\tabcolsep}{2.5pt}
\begin{tabular}{lcccccc}
\toprule
\textbf{Pathology} & \textbf{Pos.} & \textbf{Thresh.} & \multicolumn{2}{c}{\textbf{AUROC}} & \multicolumn{2}{c}{\textbf{F1-Score}} \\
\cmidrule(lr){4-5} \cmidrule(lr){6-7}
 & & & \textbf{Baseline} & \textbf{X-Fusion} & \textbf{Baseline} & \textbf{X-Fusion} \\
\midrule
Atelectasis        & 1807 & 0.42 & 0.8112 & \textbf{0.9320} & 0.3873 & \textbf{0.8870} \\
Cardiomegaly       & 390  & 0.45 & 0.8862 & \textbf{0.9580} & 0.3085 & \textbf{0.9230} \\
Consolidation      & 699  & 0.38 & 0.8013 & \textbf{0.9390} & 0.2212 & \textbf{0.9010} \\
Edema              & 348  & 0.44 & 0.8918 & \textbf{0.9620} & 0.2451 & \textbf{0.9360} \\
Effusion           & 2148 & 0.46 & 0.8827 & \textbf{0.9610} & 0.5438 & \textbf{0.9300} \\
Emphysema          & 398  & 0.40 & 0.9200 & \textbf{0.9450} & 0.3987 & \textbf{0.9130} \\
Fibrosis           & 255  & 0.32 & 0.8072 & \textbf{0.9210} & 0.1642 & \textbf{0.8890} \\
Hernia             & 30   & 0.30 & 0.8510 & \textbf{0.9810} & 0.4444 & \textbf{0.9560} \\
Infiltration       & 3166 & 0.40 & 0.7225 & \textbf{0.9180} & 0.4217 & \textbf{0.8720} \\
Mass               & 787  & 0.42 & 0.8334 & \textbf{0.9250} & 0.3827 & \textbf{0.8900} \\
Nodule             & 953  & 0.38 & 0.7555 & \textbf{0.9140} & 0.2840 & \textbf{0.8730} \\
Pleural Thick.     & 553  & 0.39 & 0.7831 & \textbf{0.9300} & 0.2065 & \textbf{0.8920} \\
Pneumonia          & 209  & 0.35 & 0.7646 & \textbf{0.9540} & 0.0792 & \textbf{0.9250} \\
Pneumothorax       & 942  & 0.43 & 0.8831 & \textbf{0.9480} & 0.4429 & \textbf{0.9170} \\
\midrule
\textbf{Macro Mean}& --   & \textbf{0.39} & 0.8162 & \textbf{0.9420} & 0.3235 & \textbf{0.9070} \\
\bottomrule
\end{tabular}
\end{table}

\begin{figure*}[!t]
\centering
\includegraphics[width=0.96\textwidth,height=3.8cm,keepaspectratio]{XFusion_Figure4_Explainability.png}
\caption{Unified Multimodal Explainability Showcase: (a) Input frontal chest radiograph; (b) Radiologist ground-truth annotation highlighting right lower lobe pneumonia consolidation; (c) Chefer-CAM visual attention rollout heatmap focusing on the opacity; and (d) Kernel SHAP physiological feature attribution quantifying vital drivers ($\text{SpO}_2$, Temperature, and Respiratory Rate).}
\label{fig:explainability}
\end{figure*}

\subsection{Ablation Analysis: Modality Synergy}
To confirm that multimodal integration provides substantive clinical gains, we conducted systematic ablation experiments summarized in Table~\ref{tab:ablation} and Fig.~\ref{fig:roc_ablation}. Adding patient vitals ($M_2$) boosts Macro-AUC by +3.78\% over vision alone, while incorporating clinical text ($M_3$) yields a +7.08\% improvement. Replacing basic feature concatenation ($M_4$) with Cross-Attention Transformer Fusion and ASL ($M_5$) unlocks an additional +4.00\% gain, establishing an overall \textbf{+12.58\% improvement} over the unimodal baseline ($p < 0.001$).

\begin{table}[t]
\caption{Ablation Study Demonstrating Modality Synergy}
\label{tab:ablation}
\centering
\scriptsize
\setlength{\tabcolsep}{3pt}
\begin{tabular}{lcccc}
\toprule
\textbf{Configuration} & \textbf{Macro-AUC} & \textbf{mAP} & \textbf{Hamming $\downarrow$} & \textbf{Subset Acc.} \\
\midrule
$M_1$: Vision-Only (BiomedCLIP)   & 0.8162 & 0.7420 & 0.0840 & 0.6120 \\
$M_2$: Vision + Metadata MLP       & 0.8540 & 0.7890 & 0.0710 & 0.6780 \\
$M_3$: Vision + CXR-BERT Text      & 0.8870 & 0.8250 & 0.0580 & 0.7340 \\
$M_4$: Tri-Modal Concatenation     & 0.9020 & 0.8410 & 0.0490 & 0.7680 \\
$M_5$: \textbf{X-Fusion (Cross-Attn + ASL)} & \textbf{0.9420} & \textbf{0.8840} & \textbf{0.0320} & \textbf{0.8420} \\
\bottomrule
\end{tabular}
\end{table}

\begin{figure}[t]
\centering
\includegraphics[width=\columnwidth,height=3.4cm,keepaspectratio]{XFusion_Figure3_ROC_Ablation.png}
\caption{Multi-label discrimination and ablation metrics across configurations.}
\label{fig:roc_ablation}
\end{figure}

\subsection{Multimodal Explainability Analysis}
The Explainability Analysis in Fig.~\ref{fig:explainability} demonstrates the dual interpretability of X-Fusion:
\begin{itemize}
    \item \textbf{Visual Localization (Chefer-CAM):} Fig.~\ref{fig:explainability}(c) reveals that attention rollout specifically highlights focal airspace consolidation in the right lower lobe, achieving high spatial alignment with the radiologist's ground-truth bounding box (Pointing Game Accuracy: $92.4\%$).
    \item \textbf{Physiological Drivers (Kernel SHAP):} Fig.~\ref{fig:explainability}(d) demonstrates that acute oxygen desaturation ($\text{SpO}_2$: $-0.38$), hyperthermia (Temp: $+0.31$), and tachypnea (RR: $+0.24$) act as primary positive clinical drivers supporting active pneumonia diagnosis.
\end{itemize}

\subsection{Computational Efficiency \& Deployment}
On a clinical workstation equipped with an Intel Core i7-13700H and NVIDIA RTX 4060 GPU (8GB VRAM), end-to-end inference—encompassing image preprocessing, CXR-BERT tokenization, cross-attention forward pass, Chefer-CAM heatmap computation, and Kernel SHAP attribution—executes in \textbf{1.18 seconds} per patient study with peak VRAM $< 5.2\text{ GB}$. The system has been operationalized within a Next.js clinical workstation communicating with a FastAPI microservice.

\section{CONCLUSION AND FUTURE WORK}
In this work, we presented \textbf{X-Fusion}, a tri-modal clinical foundation framework for multi-label chest X-ray disease diagnosis and reasoning. By harmonizing radiographic images (BiomedCLIP ViT-B/16), clinical text (CXR-BERT), and structured physiological vitals (Metadata MLP) into a unified 512-dimensional shared latent space via a 2-layer cross-attention Transformer, X-Fusion overcomes the fundamental semantic gap of unimodal CADx systems. Optimized via Asymmetric Focal Loss and InfoNCE contrastive alignment, X-Fusion achieves a state-of-the-art Macro AUROC of \textbf{0.9420} across 14 thoracic pathologies on NIH ChestX-ray14. Furthermore, its dual explainability engine (Chefer-CAM visual attention rollout and Kernel SHAP physiological attribution) equips clinicians with transparent, multi-dimensional justifications.

Future work will expand X-Fusion to multi-center clinical validation across prospective intensive care cohorts, integrate direct bounding-box coordinate regression heads for fine-grained lesion localization, and evaluate three-dimensional volumetric modalities including chest computed tomography (CT) and magnetic resonance imaging (MRI).

\balance
\section*{ACKNOWLEDGMENT}
The authors express their sincere gratitude to Dr. K. Vishal Reddy for project mentorship, domain insights, and research guidance throughout the development of X-Fusion. We also acknowledge the National Institutes of Health (NIH) Clinical Center for providing the benchmark ChestX-ray14 dataset.

\begin{thebibliography}{99}
\bibitem{wang2017chestxray8}
X. Wang, Y. Peng, L. Lu, Z. Lu, M. Bagheri, and R. M. Summers, ``ChestX-ray8: Hospital-scale chest X-ray database and benchmarks on weakly-supervised classification and localization of common thorax diseases,'' in \textit{Proc. IEEE CVPR}, 2017, pp. 2097--2106.

\bibitem{he2016deep}
K. He, X. Zhang, S. Ren, and J. Sun, ``Deep residual learning for image recognition,'' in \textit{Proc. IEEE CVPR}, 2016, pp. 770--778.

\bibitem{rajpurkar2017chexnet}
P. Rajpurkar et al., ``CheXNet: Radiologist-level pneumonia detection on chest X-rays with deep learning,'' \textit{arXiv:1711.05225}, 2017.

\bibitem{irvin2019chexpert}
J. Irvin et al., ``CheXpert: A large chest radiograph dataset with uncertainty labels and expert comparison,'' in \textit{Proc. AAAI}, 2019, pp. 590--597.

\bibitem{baltruschat2019comparison}
I. M. Baltruschat, H. Nickisch, M. Grass, T. Knopp, and A. Saalbach, ``Comparison of deep learning approaches for multi-label chest X-ray classification,'' \textit{Sci. Rep.}, vol. 9, Art. 6381, 2019.

\bibitem{allaouzi2019multilabel}
I. Allaouzi and M. Ben Ahmed, ``A novel approach for multi-label chest X-ray classification of common thorax diseases,'' \textit{IEEE Access}, vol. 7, pp. 64279--64288, 2019.

\bibitem{pham2021interpreting}
H. Pham, T. Le, D. Q. Tran, D. Ngo, and H. Nguyen, ``Interpreting chest X-rays via CNNs that exploit hierarchical disease dependencies,'' \textit{Neurocomputing}, vol. 437, pp. 186--194, 2021.

\bibitem{holste2024cxrlt}
G. Holste et al., ``Towards long-tailed, multi-label disease classification from chest X-ray: Overview of the CXR-LT challenge,'' \textit{Med. Image Anal.}, vol. 97, Art. 103224, 2024.

\bibitem{ridnik2021asymmetric}
T. Ridnik, E. Ben-Baruch, A. Noy, and L. Zelnik-Manor, ``Asymmetric loss for multi-label classification,'' in \textit{Proc. IEEE/CVF ICCV}, 2021, pp. 3343--3352.

\bibitem{dosovitskiy2021an}
A. Dosovitskiy et al., ``An image is worth 16x16 words: Transformers for image recognition at scale,'' in \textit{Proc. ICLR}, 2021.

\bibitem{liu2021swin}
Z. Liu et al., ``Swin Transformer: Hierarchical vision transformer using shifted windows,'' in \textit{Proc. IEEE/CVF ICCV}, 2021, pp. 10012--10022.

\bibitem{taslimi2022swinchex}
S. Taslimi et al., ``SwinCheX: Multi-label classification on chest X-ray images with transformers,'' \textit{arXiv:2206.04246}, 2022.

\bibitem{fu2025davit}
M. Fu, C. Tantithamthavorn, and T. Le, ``DAViT: A domain-adapted vision transformer for automated pneumonia detection and explanation,'' \textit{IEEE Access}, vol. 13, pp. 103033--103044, 2025.

\bibitem{zhang2023biomedclip}
S. Zhang et al., ``BiomedCLIP: A multimodal biomedical foundation model pretrained on fifteen million scientific image-text pairs,'' \textit{arXiv:2303.00915}, 2023.

\bibitem{boecking2022making}
B. Boecking et al., ``Making the most of text semantics to improve biomedical vision-language models,'' in \textit{ECCV}, LNCS, vol. 13696, 2022, pp. 1--21.

\bibitem{sloan2026camchex}
P. Sloan, E. Simpson, and M. Mirmehdi, ``Clinically-aligned multi-modal chest X-ray classification,'' in \textit{Proc. ML4H}, PMLR, vol. 297, 2026, pp. 228--242.

\bibitem{wang2022medclip}
Z. Wang, Z. Wu, D. Agarwal, and J. Sun, ``MedCLIP: Contrastive learning from unpaired medical images and texts,'' in \textit{Proc. EMNLP}, 2022, pp. 3876--3887.

\bibitem{aksoy2024beyond}
N. Aksoy et al., ``Beyond images: An integrative multi-modal approach to chest X-ray report generation,'' \textit{Front. Radiol.}, vol. 4, Art. 1339612, 2024.

\bibitem{singh2025chestxtranscribe}
P. Singh and S. Singh, ``ChestX-Transcribe: A multimodal transformer for automated radiology report generation,'' \textit{Front. Digit. Health}, vol. 7, Art. 1535168, 2025.

\bibitem{chefer2021transformer}
H. Chefer, S. Gur, and L. Wolf, ``Transformer interpretability beyond attention visualization,'' in \textit{Proc. IEEE/CVF CVPR}, 2021, pp. 782--791.

\bibitem{lundberg2017unified}
S. M. Lundberg and S.-I. Lee, ``A unified approach to interpreting model predictions,'' in \textit{NeurIPS}, vol. 30, 2017, pp. 4765--4774.

\end{thebibliography}

\end{document}
